\addcontentsline{toc}{chapter}{Abstract}\vspace{-1cm}
%Border
\begin{tikzpicture}[remember picture, overlay]
  \draw[line width = 4pt] ($(current page.north west) + (0.75in,-0.75in)$) rectangle ($(current page.south east) + (-0.75in,0.75in)$);
\end{tikzpicture}

Secure communication over untrusted networks is a critical challenge in modern information systems. Traditional encryption techniques rely on mathematical hardness assumptions that may be threatened by quantum computing. This project addresses the gap between classical socket communication and quantum-secure key distribution by implementing a Peer-to-Peer (P2P) messaging system using raw TCP socket programming and a quantum-inspired BB84 key exchange protocol.

The system enables users on a Local Area Network (LAN) to exchange messages securely using AES-256 encryption. The encryption keys are dynamically derived from a simulated BB84 Quantum Key Distribution (QKD) process. By independently generating and measuring qubits, Alice and Bob can compute a quantum error rate. If an eavesdropper or Man-in-the-Middle (MITM) attacker intercepts the transmission, the interference disturbs the quantum state, raising the error rate above a strict threshold. This allows the system to immediately detect passive and active attacks, rejecting compromised keys and discarding replay attempts.
A React-based dashboard provides real-time visibility into the protocol, displaying peer discovery, BB84 basis tables, encryption flows, and security alerts. The project serves as a functional demonstration of how quantum cryptographic principles can secure conventional networks against advanced threats.
